\chapter{Introduzione}

\section{Il judo e la sua storia}

Il judo nasce in Giappone alla fine dell'Ottocento per opera di Jigoro Kano, che nel 1882 fonda a Tokyo il Kodokan, il dojo da cui il judo si è diffuso in tutto il mondo~\cite{kano1986}. Kano aveva studiato per anni il jujutsu, cioè l'insieme delle tecniche di difesa a mani nude dei samurai, e ne aveva colto sia l'efficacia sia i limiti: molte tecniche erano pericolose e non potevano essere praticate liberamente senza rischiare gravi infortuni. Da questa esperienza nasce l'idea di un metodo educativo nuovo, che conservasse i principi tecnici, eliminasse le tecniche troppo pericolose e permettesse di allenarsi con un compagno per raggiungere una mutua crescita. \par

Il nome stesso riassume questa idea, infatti la parola judo significa ``via dell'adattabilità'' o ``via della flessibilità''. Invece di opporsi con la forza alla forza dell'avversario, il judoka cerca di sfruttarla, sbilanciandolo e usando il suo stesso movimento per proiettarlo o controllarlo a terra. Kano riassume il judo in due principi: \textit{Seiryoku-zenyo}, il miglior impiego dell'energia, e \textit{Jita-kyoei}, il principio della mutua crescita, secondo cui lo sviluppo del singolo passa attraverso la collaborazione con gli altri~\cite{kano1986}. Il judo nasce quindi non soltanto come arte di combattimento, ma come metodo di educazione fisica e morale. \par

\begin{figure}[ht]
    \centering
    \includegraphics[width = 0.5\textwidth]{pictures/kano.jpg}
    \creditcaption[Statua di Jigoro Kano]{Statua di Jigoro Kano, di F. Asakura, davanti al Kodokan di Tokyo.}{Foto di H. Probell, Wikimedia Commons, licenza CC BY-SA 2.5, \url{https://commons.wikimedia.org/wiki/File:Kodokan_Jigoro_Kano_Statue.jpg}}
\end{figure}

Per rendere possibili le proiezioni in sicurezza, il judo dà grande importanza alle tecniche di caduta (ukemi), che ogni praticante impara fin dalle prime lezioni e che permettono di essere proiettati senza farsi male. L'allenamento si svolge sempre in coppia sul tatami (un tappeto progettato apposta per attutire meglio le cadute) e alterna diversi tipi di esercizio: la preparazione fisica, lo studio delle tecniche, la loro ripetizione con un compagno e il combattimento libero in piedi e a terra (randori). Proprio il randori, in cui due atleti combattono liberamente cercando di applicare le proprie tecniche, è il cuore dell'allenamento e allo stesso tempo il momento in cui il corpo è sottoposto al maggior rischio di infortunio.

\section{Dal judo educativo al judo sportivo}

Nel corso del Novecento il judo si diffonde rapidamente fuori dal Giappone e assume sempre più il carattere di sport. Nel 1951 nasce la Federazione Internazionale di Judo, che unifica i regolamenti e organizza le competizioni internazionali, e nel 1964, ai Giochi Olimpici di Tokyo, il judo maschile entra per la prima volta nel programma olimpico. Dal 1992, con i Giochi di Barcellona, anche il judo femminile fa parte del programma olimpico~\cite{ijfhistory}. \par

L'ingresso del judo nel panorama olimpico ha trasformato profondamente questa disciplina. Le competizioni sono state suddivise in categorie di peso, i regolamenti sono stati più volte rivisti con l'obiettivo di rendere gli incontri più dinamici e spettacolari, e il numero di gare a livello internazionale è aumentato. Per un atleta di alto livello, il judo è diventato una disciplina che richiede, oltre alla tecnica, una preparazione fisica molto sviluppata. Gli studi sul profilo fisiologico degli atleti d'élite evidenziano che il judo richiede soprattutto un'elevata potenza anaerobica, in particolare della parte superiore del corpo, una buona capacità aerobica e una bassa percentuale di massa grassa, con l'eccezione delle categorie più pesanti~\cite{franchini2011}. \par

\begin{figure}[ht]
    \centering
    \includegraphics[width = 0.8\textwidth]{pictures/throw.jpg}
    \creditcaption[Proiezione di judo]{Finale pesi massimi, Parigi 2024.}{Foto \textcopyright~International Judo Federation, \url{https://www.ijf.org/competition/2653/photos?id_weight=7}}
\end{figure}

Di conseguenza, anche l'allenamento ha subito cambiamenti. Un atleta di alto livello si allena spesso più volte al giorno, alternando sedute tecniche, sedute di randori ad alta intensità e sedute di preparazione fisica, e affronta periodi di preparazione intensa in vista delle gare più importanti. Nei raduni e nei centri di allenamento nazionali capita spesso di trovare molti atleti sullo stesso tatami, che si allenano a coppie cambiando frequentemente compagno. A queste esigenze si aggiunge, per molti atleti, la necessità di rientrare nella propria categoria di peso, che comporta talvolta perdite di peso rapide nei giorni che precedono la gara. \par

Questa evoluzione ha portato a prestazioni sempre più elevate, ma ha anche aumentato le sollecitazioni a cui è sottoposto il corpo degli atleti. Un carico di allenamento elevato e mal distribuito, unito a un recupero insufficiente, è oggi considerato uno dei principali fattori di rischio di infortunio nello sport~\cite{soligard2016}. Non è però il carico in sé a essere dannoso, infatti un carico elevato, raggiunto in modo graduale, rende gli atleti più resistenti, mentre sono gli aumenti bruschi e non accompagnati da un adeguato recupero a far crescere il rischio~\cite{gabbett2016}.

\section{Gli infortuni nel judo}

Il judo è uno sport di combattimento e di contatto, nel quale gli atleti vengono continuamente proiettati, cadono, si oppongono con forza alle tecniche dell'avversario e lottano per ottenere il controllo a terra. Non sorprende quindi che gli infortuni siano frequenti. Una revisione sistematica degli studi sugli infortuni nel judo indica che i più comuni sono le distorsioni, gli stiramenti e le contusioni, che interessano soprattutto il ginocchio, la spalla e le dita, e che il meccanismo più frequente è l'essere proiettati. Gli infortuni gravi, come quelli alla testa e alla colonna vertebrale, sono rari, mentre tra i problemi cronici si riscontrano soprattutto quelli alle articolazioni delle dita, alla parte bassa della schiena, alle anche e alle ginocchia~\cite{pocecco2013}. \par

\begin{figure}[ht]
    \centering
    \includegraphics[width = 0.8\textwidth]{pictures/training-camp.jpg}
    \creditcaption[Tatami durante una sessione di allenamento]{Tatami durante una sessione di allenamento, Judo Training Camp, Roma 2026.}{Foto \textcopyright~FIJLKAM, \url{https://www.fijlkam.it/media-fijlkam/galleria-fotografica/galleria-fotografica/2026/judo-training-camp-roma-2026.html}}
\end{figure}

Ad alto livello il judo risulta tra gli sport di combattimento più rischiosi. Un'analisi degli infortuni negli sport di combattimento in tre Olimpiadi consecutive, Pechino 2008, Londra 2012 e Rio de Janeiro 2016, ha mostrato che il judo ha il tasso di infortuni più alto tra gli sport considerati, con circa $9.6$ infortuni ogni $1000$ minuti di gara. Inoltre, più di un terzo degli infortuni nel judo, il $35.9\%$, ha costretto l'atleta a fermarsi per più di una settimana~\cite{lystad2021}. \par

Gli infortuni, però, non si verificano soltanto in gara. La maggior parte del tempo di un atleta è dedicata all'allenamento, ed è proprio lì che si accumulano la fatica e le sollecitazioni ripetute. Uno studio prospettico condotto per quattro anni nel centro di allenamento della nazionale coreana di judo ha registrato $782$ infortuni, pari a circa quattro infortuni per atleta ogni anno. Circa la metà erano lievi, con meno di quattro giorni di cure, e la zona più colpita era la parte inferiore del corpo, con il $44.2\%$ degli infortuni, seguita dalla parte superiore, con il $29.8\%$, e dal tronco, con il $20.3\%$~\cite{kim2015}. \par

Oltre al carico di allenamento, la letteratura indica altri fattori che possono aumentare il rischio. Uno di questi è la perdita di peso prima delle gare: uno studio sugli infortuni durante le competizioni ha osservato che gli atleti che avevano perso peso prima della gara si infortunavano più spesso, e gli autori sconsigliano questa pratica~\cite{green2007}. Altri fattori discussi sono l'alimentazione, l'idratazione e gli aspetti psicologici~\cite{pocecco2013}. \par

Le conseguenze di un infortunio vanno oltre il dolore e le cure. Un atleta infortunato deve interrompere o ridurre l'allenamento, perde parte della preparazione e, se torna ad allenarsi prima di essere completamente guarito, può essere più esposto a nuovi infortuni. Per una squadra o una palestra, inoltre, gli infortuni riducono il numero di atleti disponibili e possono portare alcuni di loro ad abbandonare l'attività.

\section{Perché un modello matematico}

Gli studi citati descrivono con precisione quanti infortuni avvengono, di che tipo sono e in quali parti del corpo si localizzano. Ciò che risulta più difficile da studiare è invece il modo in cui gli infortuni si sviluppano nel tempo all'interno di un gruppo che si allena insieme. In una palestra o in un raduno, infatti, gli infortuni non sono eventi isolati: la fatica si accumula seduta dopo seduta e aumenta il rischio per tutti gli atleti; un tatami affollato rende più probabili gli scontri tra coppie; un atleta che ha subito un infortunio lieve continua spesso ad allenarsi in condizioni non ottimali ed è quindi più esposto a un nuovo infortunio; e gli atleti che si infortunano sono anche quelli che più facilmente abbandonano il gruppo. \par

Questi fenomeni sono legati tra loro, e il loro effetto complessivo non è facile da prevedere in modo intuitivo. Ridurre il carico di una seduta, ad esempio, diminuisce la fatica e quindi gli infortuni, ma riduce anche il beneficio dell'allenamento; aumentare la prevenzione riduce gli infortuni, ma richiede tempo e attenzione; e una riduzione degli infortuni ha a sua volta effetti sulla composizione del gruppo e sugli abbandoni. Studiare queste interazioni direttamente sugli atleti sarebbe difficile e, in molti casi, non accettabile: non è possibile sottoporre un gruppo a un carico eccessivo solo per osservare quanti infortuni produce. \par

Un modello matematico permette invece di studiare queste interazioni in modo controllato. Descrivendo con delle equazioni come cambia il gruppo da una seduta all'altra, è possibile simulare periodi di allenamento diversi, confrontare scenari differenti e valutare l'effetto di una strategia di prevenzione prima di applicarla. Il modello, naturalmente, è una semplificazione della realtà, e i suoi risultati vanno interpretati con cautela: non sostituisce l'esperienza dell'allenatore né i dati raccolti sul campo, ma offre uno strumento per ragionare in modo quantitativo su fenomeni che altrimenti si possono valutare solo a intuito. \par

In questa tesi proponiamo un modello di questo tipo. Il modello divide gli atleti di un gruppo in quattro popolazioni, cioè i sani, gli atleti con infortuni pregressi, gli infortunati e gli atleti che si allenano nonostante un infortunio lieve, e descrive come gli atleti passano da una popolazione all'altra in funzione della fatica accumulata, della densità di atleti sul tatami e dei flussi di ingresso e di abbandono del gruppo. Al modello viene poi aggiunto un controllo, che rappresenta la prevenzione messa in atto dall'allenatore, per studiare quanto e come la prevenzione possa ridurre gli infortuni.